\section{Reinforcement Learning}
\label{sec:rl}

\gls{rl} is the branch of \gls{ml} designed to learn an optimal policy in a sequential decision-making problem, modelled by an \gls{mdp}. 
The main difficulty of \gls{rl} is to collect data by interacting with the environment and learning from these data.
The \emph{sample complexity}--the number of samples an agent collects from interacting with its environment--is an important metric for \gls{rl} algorithms. 
Sampling from an environment can be costly in terms of time, computational resources, or money. 
Simulations could be run instead, but this introduces problems when translating the model from simulation to reality\footnote{Notably, the delay is typically ignored in simulation and can pose problems when testing in the real world.}. 
Therefore, it is common to use sample complexity as a tool to measure the efficiency of a \gls{rl} algorithm. 

Although the performance of an agent is still the final goal that the designer would like to optimise, there can be some trade-off between the final goal and the sample complexity, as a fast learning agent may be preferred over an agent guaranteed to reach an optimal behaviour but in a proscriptively long time. 
This trade-off is known as the exploration-exploitation problem, where the agent must trade off between taking actions whose outcomes he knows to be beneficial against taking actions whose outcomes he is not certain about, in order to gain knowledge on its environment. 

Historically, the first methods applied to \gls{mdp} problems were \emph{dynamic programming} algorithms such as \emph{policy iteration} \cite{howard1960dynamic}. 
But these methods are limited because they require knowledge of the transition dynamics as well as maintaining a memory of the value function in each state of the \gls{mdp}.
Moving toward more realistic problems, \gls{rl} has removed the hypothesis of known dynamics and has started considering larger state spaces, eventually continuous ones, which therefore require function approximation. 
We will dive into some of these solutions in the remainder of this section. 
They will be useful in this dissertation, as many algorithms in delayed \gls{rl} are based upon them.
We will first present algorithms designed for the discounted return criteria and finish with an algorithm for the average reward one.

\subsection{Value-Based}
The first approach to solving \gls{rl} problems is to learn the value function $V$ or the action value function $Q$, which is what value-based approaches do. 
Seminal approaches are SARSA \cite{rummery1994line} and Q-learning \cite{watkins1989learning}, which are based on dynamic programming.

\subsubsection{Q-Learning}
Q-learning \cite{watkins1989learning} is a popular algorithm that recursively learns the optimal Q function by interacting with the environment and updating its estimate via the optimal Bellman \cref{eq:bellman_optimal_equation_Q}.
The algorithm samples tuples of the form $(s_t,a_t,r_t,s_{t+1})$ by interacting with the environment and uses them to update the Q-function as follows,
\begin{align}
\label{eq:update_q_learning}
    \hat Q(s_t, a_t)_i = \hat Q(s_t, a_t) + \alpha \delta_{TD,t},
\end{align}
where $\delta_{TD,k}$ is called the \gls{td} error between the current Q-function and its one-step bootstrap. Formally, at step $k$, it reads:
\begin{align}
\label{eq:td_error_q_learning}
    \delta_{TD,t} = r_t+\gamma \max_{a\in\actionspace} \hat Q(s_{t+1}, a) -\hat Q(s_t, a_t).
\end{align}
Q-learning is an \emph{off-policy} algorithm because it uses a different policy to sample from the environment than the one it uses for its update. 
Sampling is usually made with an $\epsilon$-greedy policy with respect to $\hat{Q}$, which means that with probability $1-\epsilon$ this policy takes the action that maximises $\hat{Q}$ and with probability $\epsilon$ selects an action uniformly at random in $\actionspace$. Instead, the policy used in the \gls{td} error term of the update (\Cref{eq:td_error_q_learning}) is greedy with respect to $\hat Q$ as it always selects the maximising action. 
The sampling $\epsilon$-greedy policy (or variations) is useful for exploration purposes. 
It ensures that each state-action pair of the \gls{mdp} will be visited infinitely often, which is a requirement for the proof of the convergence of the algorithm \cite{watkins1992q}.
It allows the agent to assess the effects of different actions, potentially yielding higher rewards.

\subsubsection{SARSA}
The SARSA \cite{rummery1994line} algorithm is very similar to Q-learning. 
It is an \emph{on-policy} algorithm; the agent uses its current policy to sample tuples $(s_t,a_t,r_t,s_{t+1},a_{t+1})$ from the environment.
Using these samples, the \gls{td} error expression now reads:
\begin{align*}
    \delta_{TD,t} = r_t+\gamma\hat Q(s_{t+1}, a_{t+1})-\hat Q(s_t, a_t).
\end{align*}
Note that $a_{t+1}$ replaces the action that maximises $Q$ in the update of \Cref{eq:td_error_q_learning}.

\subsubsection{Deep Q network}
When the state space $\statespace$ grows too large, a solution is to resort to approximation for the Q-function. 
A choice for such approximators is \emph{deep neural network} as in \gls{dqn} \cite{mnih2013playing}.
Note that the action space $\actionspace$ remains discrete, so the maximum of \Cref{eq:td_error_q_learning} can be computed.
Relaxing the assumption of a discrete action space has been proposed by other approaches, such as Deep Deterministic Policy Gradient (DDPG) \cite{lillicrap2015continuous}.

In \gls{dqn}, the Q-function is now parameterized by some parameter $\boldsymbol\theta$ with the goal of finding $Q_{\boldsymbol\theta}(s,a) \sim Q^\star(s,a)$. 
Given the current set of parameters $\boldsymbol\theta_i$, the next set of parameters $\boldsymbol\theta_{i+1}$ is obtained by minimising the following \gls{td} error:
\begin{align}
\label{eq:update_dqn}
    \delta_{TD}(\theta_i,\rho) = \expectedvalue_{(s,a,r,s')\sim\rho}\left[\left(r+\gamma\max_{a\in\actionspace}Q_{\boldsymbol\theta_{i}}(s') - Q_{\boldsymbol\theta_{i+1}}(s,a)\right)^2\right].
\end{align}
where $\rho$ is a distribution over $\statespace\times\actionspace\times[0,\maximumreward]\times\statespace$. 
Usually, at each iteration of the minimisation of \Cref{eq:update_dqn}, new samples $(s,a,r,s')$ are collected with an $\epsilon$-greedy policy \wrt $Q_{\boldsymbol\theta_{i+1}}$. 
\cite{mnih2013playing} observed that keeping not only the last sampled tuples but also the older samples in a replay buffer to minimise \Cref{eq:update_dqn} smooths learning and makes the problem more similar to \gls{sl}. This technique is known as \emph{experience replay}.
The buffer is therefore filled with samples from many different policies, and we call this mixture of policies $\rho$.

\subsection{Policy-Based}
\label{subsec:policy_based}
One main limitation of value-based approaches is that the agent learns extra knowledge about the environment that might not be of direct usefulness. 
In the end, one is only interested in learning optimal behaviour, and learning the Q-function is only an indirect way to reach this goal. 
Instead, policy-based approaches learn the policy directly from interactions with the environment.

In general, policies are considered a parameterised function to shift the problem from searching over a set of functions to searching over a set of parameters.
This framework, called \gls{po} \cite{deisenroth2013survey} considers the policy in a parametric set $\Pi_{\Theta} = \{\pi_{\vectorialform{\theta}} : \vectorialform{\theta} \in \Theta \subseteq \realnumbers^{n_\theta}\}$.

The policy is entirely specified by its parameter $\vectorialform{\theta}$ and, therefore, we simplify the notation for its return as $\expectedreturn[][](\vectorialform{\theta})=\expectedreturn[\gamma][\pi_\vectorialform{\theta}]$.
The new objective of the agent then becomes,
\begin{align}
\label{eq:objective_policy_opt}
	\vectorialform{\theta}^\star \in \argmax_{\vectorialform{\theta} \in \Theta} \expectedreturn[][](\vectorialform{\theta}).
\end{align}

The policy $\pi_{\vectorialform{\theta}}$ needs to be stochastic in order to provide enough exploration. 
If the policy is also differentiable in $\vectorialform{\theta}$, then the following fundamental result of policy optimisation holds.

\begin{thm}[Policy Gradient Theorem]
    Let $\pi_{\vectorialform{\theta}}\in\Pi_{\Theta}$ be a stochastic policy, differentiable in $\vectorialform{\theta}$, then the gradient of its return reads
    \begin{align*}
        \nabla_{\vectorialform{\theta}} \expectedreturn[][](\vectorialform{\theta}) =  \frac{1}{1-\gamma}\expectedvalue_{\substack{s\sim\discountedstateoccupancydistribution[\gamma][\pi_{\vectorialform{\theta}}]\\a\sim\pi_{\vectorialform{\theta}}(\cdot\vert s)}}\left[Q^{\pi_{\vectorialform{\theta}}}(s,a)\nabla_{\vectorialform{\theta}} \log \pi_{\vectorialform{\theta}}(a\vert s )\right]
    \end{align*}
\end{thm}
Note that in many works, including state-of-the-art methods, an approximation of the gradient is used, dropping the discounting included in the state distribution but retaining it in the Q-function. 
This quantity may no longer be the gradient of any function \cite{nota2020policy}, yet it provides good empirical results. 

The policy gradient does not require access to the \gls{mdp} model and can be estimated from samples. 
We present two of the main estimators in the following paragraphs. 
Their formulations assume access to a dataset of trajectories $(\tau_i)_{1\le i\le n}$ where each trajectory is of length $H$, specifically $\tau_i=(s_0^i,a_0^i,r_0^i\dots,s_{H}^i)$.

\subsubsection{REINFORCE} 
Introduced by \cite{williams1992simple}, the REINFORCE estimator reads,
\begin{align*}
    \widehat{\nabla_{\vectorialform{\theta}} \expectedreturn[][](\vectorialform{\theta})} = \frac{1}{n}\sum_{i=1}^{n}\sum_{t=1}^{H-1} \nabla_{\vectorialform{\theta}} \log \pi_{\vectorialform{\theta}}(a_t^i\vert s_t^i)\left(\sum_{t'=0}^{H-1}\gamma^{t'} r_{t'}^i\right).
\end{align*}
This estimator usually has a prohibitively high variance.
However, a constant baseline $b$ can be subtracted from the rightmost term without added bias while allowing control of the variance 
\cite{williams1992simple}.

\subsubsection{Policy Gradient Theorem} 
Another approach to reducing the variance of the previous estimator is obtained by noting that future actions do not influence past rewards. 
Indeed, for any $i>j$, note $\discountedstateoccupancydistribution[\gamma,j:i][\pi]$ the distribution of the chunk of trajectory $(s_j,a_j,\dots,s_i,a_{i})$ under policy $\pi$, then one has
\begin{align*}
    &\expectedvalue_{(s_j,a_j,\dots,s_i,a_{i})\sim\discountedstateoccupancydistribution[\gamma,j:i][\pi_{\vectorialform{\theta}}]}\left[\nabla_{\vectorialform{\theta}} \log \pi_{\vectorialform{\theta}}(a_i\vert s_i)r(s_{j},a_{j})\right] 
    \\&\quad= \expectedvalue_{\substack{(s_j,a_j,\dots,s_{i-1},a_{i-1})\sim\discountedstateoccupancydistribution[\gamma,j:i-1][\pi_{\vectorialform{\theta}}]\\s_i\sim p(\cdot\vert s_{i-1},a_{i-1})}}\left[ r(s_{j},a_{j})\int_{\actionspace}\pi_{\vectorialform{\theta}}(a_i) \nabla_{\vectorialform{\theta}} \log \pi_{\vectorialform{\theta}}(a_i\vert s_i)\;\de a_i\;\right]
    \\&\quad= \expectedvalue_{\substack{(s_j,a_j,\dots,s_{i-1},a_{i-1})\sim\discountedstateoccupancydistribution[\gamma,j:i-1][\pi_{\vectorialform{\theta}}]\\s_i\sim p(\cdot\vert s_{i-1},a_{i-1})}}\left[r(s_{j},a_{j})\int_{\actionspace}\nabla_{\vectorialform{\theta}}\pi_{\vectorialform{\theta}}(a_i) \;\de a_i\; \right]
    \\&\quad= \expectedvalue_{\substack{(s_j,a_j,\dots,s_{i-1},a_{i-1})\sim\discountedstateoccupancydistribution[\gamma,j:i-1][\pi_{\vectorialform{\theta}}]\\s_i\sim p(\cdot\vert s_{i-1},a_{i-1})}}\left[r(s_{j},a_{j})\nabla_{\vectorialform{\theta}}1\right]
    \\&\quad= 0
\end{align*}
This yields the Policy Gradient Theorem (PGT) \cite{sutton1999policy},
\begin{align}
\label{eq:pgt}
    \widehat{\nabla_{\vectorialform{\theta}} \expectedreturn[][](\vectorialform{\theta})} = \frac{1}{n}\sum_{i=1}^{n}\sum_{t=1}^{H-1} \nabla_{\vectorialform{\theta}} \log \pi_{\vectorialform{\theta}}(a_t^i\vert s_t^i)\left(\sum_{t'=t}^{H-1}\gamma^{t'} r_{t'}^i-b(s_t)\right).
\end{align}
Note how the sum in the rightmost term now starts at $t$. 
Here, also, a baseline $b$ has been introduced, with the aim of further reducing the variance, as in REINFORCE. 
In PGT, the baseline might depend on the current state $s_t$ without added bias \cite{peters2008reinforcement}.
Smart baseline choices are discussed in \cite{peters2008reinforcement}.


\subsubsection{Trust Region Policy Optimisation}
\gls{trpo}\cite{schulman2015trust} is a policy-based algorithm whose idea is to provide safe improvement steps, i.e. updates of the policy that guarantee improvement in performance.
The theoretical results upon which it is based rely on the famous performance difference lemma, which we recall here.
\begin{lemma}[Performance Difference Lemma, Lemma 6.1 of \cite{kakade2002approximately}]
    Let $\pi,\pi'\in\Pi^{SR}$ with respective expected discounted return $\expectedreturn[\gamma][\pi]$ and $\expectedreturn[\gamma][\pi']$ then,
    \begin{align*}
        \expectedreturn[][\pi'] - \expectedreturn[][\pi] = \frac{1}{1-\gamma}\expectedvalue_{(s,a)\sim\discountedstateoccupancydistribution[\gamma][\pi']}\left[\advantagefunction^{\pi}(s,a)\right].
    \end{align*}
\end{lemma}

Based on this lemma, \cite{schulman2015trust} build a surrogate objective whose maximisation guarantees an improvement of $\expectedreturn[\gamma][\pi']$.

Letting $\textstyle C=\frac{4\varepsilon\gamma}{(1-\gamma)^2}$ and $\varepsilon=\max_{s,a}\advantagefunction^\pi(s,a)$ \cite{schulman2015trust} show that 
\begin{align*}
    \expectedreturn[\gamma][\pi'] \ge \underbrace{\expectedreturn[\gamma][\pi] + \frac{1}{1-\gamma}\expectedvalue_{(s,a)\sim\discountedstateoccupancydistribution[\gamma]}\left[\advantagefunction^{\pi'}(s,a)\right]}_{\coloneq L^\pi(\pi')} - C \max_s \kullbackleiblerdivergence^{\max} (\pi(\cdot\vert s),\pi'(\cdot\vert s)).
\end{align*}

To simplify the optimisation, in their final objective, \cite{schulman2015trust} substitute the average KL for the maximum. 
The former is indeed simpler to compute in practice.
Furthermore, the penalisation term in $C$ is replaced by a constraint in order to allow for larger yet robust steps.
The maximisation becomes,
\begin{align*}
\begin{array}{cl}
     \underset{\pi'}{\text{maximise}} &  L^\pi(\pi')\\
     \text{subject to} & \expectedvalue_{s\sim \discountedstateoccupancydistribution[\gamma][\pi]}\left[ \kullbackleiblerdivergence (\pi(\cdot\vert s),\pi'(\cdot\vert s))\right]\le \delta.
\end{array}
\end{align*}
\gls{trpo} has proved to be an empirically efficient algorithm. 


\subsubsection{Proximal Policy Optimisation}
Despite its empirical success, \gls{trpo} has a high computational cost because it involves the computation of the Hessian of the Kullback-Leibler divergence and its inverse.
Instead, \gls{ppo} \cite{schulman2017proximal} propose to cast the problem back to a first-order optimisation scheme using a clipping operation to maintain the policy in a safe region. 
Empirical results suggest that \gls{ppo} is as efficient as \gls{trpo} if not better, particularly for its simpler computations.  


\subsubsection{Parameter-based Optimisation}
\label{subsec:param_based_optim}
From the parameterised \gls{po}, it is possible to add a layer of abstraction and consider optimising for \emph{hyper-policies} instead of policies. 
Basically, a hyper-policy defines a rule for selecting the policy that the agent will follow.
In this setting, the parameters $\vectorialform{\theta}$ of the policy are sampled by a hyper-policy $\nu_{\vectorialform{\rho}}$ which is itself parameterised by a vector $\vectorialform{\rho}$ in $\mathcal{V}_{\mathcal{P}} = \{\nu_{\vectorialform{\rho}} : \vectorialform{\rho} \in \mathcal{P} \subseteq \realnumbers^{n_\rho}\}$~\cite{sehnke2008parameter}.
Moving the stochasticity to the hyper-policy, the policy need not be stochastic anymore. 
This is a great advantage for reducing the variance of the estimators. 
Indeed, a trajectory can be collected by sampling only a single $\vectorialform{\theta}$. 
The stochasticity of the trajectory then only results from the environment.
The objective in this framework becomes,
\begin{align}
\label{eq:objective_hyperpolicy_opt}
	\vectorialform{\rho}^\star \in \argmax_{\vectorialform{\rho} \in \mathcal{P}} \expectedreturn[][](\vectorialform{\rho})=\argmax_{\vectorialform{\rho}\in \mathcal{P}} \expectedvalue_{\boldsymbol\theta\sim \nu_{\vectorialform{\rho}}}\expectedreturn[][](\vectorialform{\theta}).
\end{align}
We refer to this setting as parameter-based \gls{po} and to the original setting as action-based \gls{po}. 


\subsection{Actor-Critic}
In the formulation of \Cref{eq:pgt}, the term $\textstyle\sum\nolimits_{t'=t}^{H-1}\gamma^{t'} r_{t'}$ is no less than a Monte-Carlo estimate of the Q-function $Q(s_t,a_t)$. 
\emph{Actor-critic} methods propose to change this estimate to a \gls{td} estimate such as TD$(0)$ which we note:
\begin{align*}
    G_{t}^0 = r_{t} + \gamma V_{\boldsymbol\phi}(s_{t+1}).
\end{align*}
Here, $V_{\boldsymbol\phi}$ is the current approximation of the state value function using the parameters $\boldsymbol\phi$.
Choosing $V_{\boldsymbol\phi}(s_{t})$ as the baseline in \Cref{eq:pgt} is an interesting choice made in the Advantage Actor-Critic (A2C) algorithm \cite{mnih2016asynchronous}. 
It transforms the rightmost parentheses in \Cref{eq:pgt} into:
\begin{align}
\label{eq:a2c}
    \widehat{\nabla_{\vectorialform{\theta}} \expectedreturn[][](\vectorialform{\theta})} = \frac{1}{n}\sum_{i=1}^{n}\sum_{t=1}^{H-1} \nabla_{\vectorialform{\theta}} \log \pi_{\vectorialform{\theta}}(a_t^i\vert s_t^i)
    \left(\vphantom{\sum_{i=1}^{n}\sum_{t=1}^{H-1}} r_{t} + \gamma V_{\boldsymbol\phi}(s_{t+1}) - V_{\boldsymbol\phi}(s_{t})\right).
\end{align}
Note that the new term corresponds to $r_{t} + \gamma V_{\boldsymbol\phi}(s_{t+1}) - V_{\boldsymbol\phi}(s_{t})=\advantagefunction_{\boldsymbol\phi}(s_t,a_t)$ which is an estimate of the advantage function (see~\Cref{eq:advantage_function}).
The name actor-critic comes from the fact that the method learns both a policy $\pi_{\vectorialform{\theta}}$, the \emph{actor}, and a state value function $V_{\boldsymbol\phi}$, the \emph{critic}. 
The algorithm alternatively updates the policy with the gradient estimate of \Cref{eq:a2c} and the value function to minimise the expected squared \gls{td} error (see \Cref{eq:update_dqn} for comparison):
\begin{align}
\label{eq:update_a2c}
    \delta_{TD}(\boldsymbol\phi_i,\rho) = \expectedvalue_{(s,a,r,s')\sim\rho}\left[\left(r+\gamma V_{\boldsymbol\phi_i}(s') - V_{\boldsymbol\phi_{i}}(s)\right)^2\right].
\end{align}
Again, $\rho$ is the distribution in some dataset of transitions $\mathcal D$.
One major advantage of actor-critic over PGT and REINFORCE is that it enables off-policy learning.


\subsubsection{Soft Actor Critic}
An example of an actor-critic algorithm we will use in this thesis is the \gls{sac} \cite{haarnoja2018soft}. 
Its main feature is to add an entropy regularisation to the update of the regular A2C described above.
Entropy is defined for some random variable $X$ with probability $p(x)=\probability(X=x)$ as
\begin{align*}
    \entropy(X) = \expectedvalue_{x\sim p}[-\log p(x)].
\end{align*}
Entropy-regularised \gls{rl} adds a bonus to the entropy of the policy to the original reward function. 
This implies that the return now reads,
\begin{align}
\label{eq:discounted_reward_entropy}
    \expectedreturn[\gamma] = (1-\gamma)\expectedvalue\left[\sum_{t=1}^H \gamma^t  \probability\left(S_t=s; \pi,\mu\right)\pi(a\vert s) \left(r(s,a)v + \alpha \entropy(\pi(\cdot\vert s))\right)\right],
\end{align}
with $\alpha>0$ a parameter that controls the regularisation.

From this new return, a new state value function and state-action value function can be proposed. The \gls{td} error changes from \Cref{eq:update_a2c} as well. 
Therefore, in \gls{sac}, the state-action value function is learnt by minimising 
\begin{align}
\label{eq:update_sac}
    \delta_{TD}(\boldsymbol\phi_i,\rho) &= \quad\expectedvalue_{(s,a,r,s')\sim\rho}\left[\left(y(r,s') - Q_{\boldsymbol\phi_{i}}(s,a)\right)^2\right],
\end{align}
where
\begin{align*}
    y(r,s') = \expectedvalue_{a'\sim\pi_{\boldsymbol\theta_{i}}}\left[r+\gamma \min_{j\in\{1,2\}}Q_{\boldsymbol\phi_{\text{target,j}}}(s',a')-\alpha \log\pi_{\boldsymbol\theta_{i}}(a'\vert s') \right].
\end{align*}
The target Q-functions $(Q_{\boldsymbol\phi_{\text{target,j}}})_{j\in\{1,2\}}$ are Polyak averages of the parameters of the Q-functions over the last iterations. They are used instead of the current Q-function $Q_{\boldsymbol\phi_i}$ to avoid the instability that is empirically observed otherwise.
The $\textstyle\min_{j\in\{1,2\}}$ is another empirical trick to avoid overestimation. 
We refer to \cite{haarnoja2018soft} for the full algorithm.


\subsection{Imitation Learning}
\label{subsec:imitation_learning}
\emph{Imitation learning} is a paradigm that comes from the observation that it is usually easier to learn a skill from demonstrations than learning it from scratch. 
In fact, learning to imitate the policy of an expert on a fixed set of observations falls under the \emph{\glsfirst{sl}} framework, which is usually simpler than \gls{rl}.
Imitation learning, therefore, aims to close the gap between \gls{rl} and \gls{sl}.
To define the framework more formally, let $\pi_E$ be an expert policy; most imitation learning approaches aim to find the policy of the learner, $\pi_I$, which minimises $\textstyle\expectedvalue_{s\sim d_\mu^{\pi_E}}[l(s,\pi_I)]$ \cite{ross2011reduction}. 
The function $l(s,\pi)$ is a loss designed to make $\pi_I$ more similar to $\pi_E$.
Notably, this objective is defined under the state distribution induced by $\pi_E$.
This can be problematic because, whenever the learner makes an error, it may end up in a state where its knowledge of the expert's behaviour is poor and make further errors.
Thus, errors can propagate as the squared effective horizon \cite[Theorem~2.1]{ross2010efficient}\cite[Theorem~1]{xu2020error}.

Other imitation learning approaches are designed to solve this problem, but are usually not practical, as explained by \cite{ross2011reduction}. 
For instance, \cite{ross2010efficient} propose the use of a non-stationary policy, which requires training a policy for each time step, which is computationally expensive for long horizons. 
Another approach by \cite{ross2010efficient}--SMILe--considers a mixture of policies augmented by one policy at each iteration. 
As stated by \cite{ross2011reduction}, this is problematic in practice since the mixture contains policies of different qualities. 
This could clearly create instability, as weaker policies are likely to make errors, while stronger policies must compensate for these errors.

A rather successful solution to the aforementioned problem is called Dataset Aggregation (\textsc{DAgger})~\cite{ross2011reduction}. 
This imitation algorithm computes its loss under the learner's state distribution \cite{osa2018algorithmic} and is, therefore, able to account for the shift in distribution induced by the learner's policy-making potential errors. 
As we shall later see, this is a sought-after property for an imitation algorithm in the case of delays since the delay can exacerbate the difference in state distribution between a delayed and an undelayed policy.
Practically, the idea of \textsc{DAgger} is simple; new samples are collected according to a policy similar to that of the learner, and the expert is only queried on those samples afterwards to understand what it would have done instead.
This has the great advantage of providing samples that match the learner's state distribution.  
The sampling policy is $\pi_i = \beta_i \pi_E + (1-\beta_i) \hat{\pi}_i$, where $\beta_i$ is the weight of a mixture of the expert policy and the current policy of the learner $\hat{\pi}_i$. 
A dataset $\mathcal{D}$ for the \gls{sl} step is built by adding the sampled states $s$ with the action $\pi_E(s)$ selected by the expert. 
Then, a new imitated policy $\hat{\pi}_{i+1}$ is trained on $\mathcal{D}$. 
The sequence $(\beta_i)_{i\in[\![1,N]\!]}$ is such that $\beta_1=1$, so as to sample initially only from $\pi_E$ and $\beta_N=0$ to sample only from the imitated policy in the end. 


\subsection{Theoretical Reinforcement Learning}
\label{subsec:theoretical_rl}

In this last subsection, we will introduce approaches that are designed for the average reward criteria.
These approaches are usually more theoretical.
They are interested in balancing exploration of the \gls{mdp} and exploitation of acquired knowledge.
The agent must learn to trade-off between these strategies so as not to lose too much with respect to an expert.
Losing can be intended in terms of missed opportunities for lack of exploration or missed rewards for lack of exploitation.
This setting belongs to \emph{online learning}.
To better understand the setting, let us present a subfield of online learning, that is \emph{prediction of individual sequences} \cite{cesa2006prediction}.
It considers an agent who has to choose repeatedly between $K$ actions, called arms. 
Each arm provides the agent with a reward, which can either be generated in a stochastic or adversarial manner. 
The environment has no state and the agent can always choose between the $K$ arms at any time. 
It is possible to represent this framework as an \gls{mdp} with a unique state, \ie $\cardinal{\statespace}=1$. 
When the agent can only observe the reward of the arm it has chosen and not the reward of the $K-1$ other arms, the setting is called a \emph{multi-armed bandit} \cite{lattimore2020bandit}. When the reward is stochastic, the bandit is called \emph{stochastic bandit} while if it is adversarial, it is referred to as \emph{adversarial bandit}.
In this field, an important concept is regret, that is, the expected difference between the reward collected by the learning agent and by the best policy in hindsight. 
The goal is to get a sub-linear regret, indicating that the agent is getting closer and closer to the expert's policy.
A common approach to this problem is to use \emph{optimism} in the face of uncertainty. 
An upper-confidence bound is computed on the estimated quantities, the reward of the arms, and the agent selects the arm with the best upper bound. We shall see in \Cref{subsec:reward_delay_related} how the delay is considered in the bandit literature.

The idea of upper-confidence bounds has also been applied to \gls{rl}. 
A precursor was the UCRL algorithm \cite{auer2006logarithmic} where confidence bounds are maintained on the reward and transition functions.
The regret in \gls{rl}, for some \gls{mdp} $\markovdecisionprocess$ is defined as:
\begin{align*}
    \mathcal{R}(T,s) \coloneq T\cdot\max_{\pi\in\Pi^{\text{MR}}} \averagereturnfunction(s) - \sum_{t=1}^T \expectedrewardfunction(s_t,a_t)
\end{align*}
where $\averagereturnfunction(s)$ is the average reward performance of the best Markovian policy with initial state distribution $\delta_s$ and $\textstyle \sum_{t=1}^T \expectedrewardfunction(s_t,a_t)$ is the reward collected by the learning algorithm during the first $T$ steps.
UCRL uses information on the mixing time of the \gls{mdp} and assumes the ergodicity of the process in order to achieve sub-linear regret. 

As we will see in the following example, later approaches have enhanced the results of UCRL in terms of smaller regret and weaker assumptions.
From these approaches, we focus on UCRL2~\cite{auer2008near}, which applies to \emph{communicating} \glspl{mdp} \cite[Section~8.3.1]{puterman1994markov}.  
We recall the definition of communicating \glspl{mdp} below.

\begin{defi}[Communicating \gls{mdp}]
\label{def:communicating_mdp}
    An \gls{mdp} $\markovdecisionprocess$ is said to be communicating if, for any states $s$ and $s'$, there exists a deterministic and stationary policy that has a non-zero probability of reaching $s'$ starting from $s$.
\end{defi}

Note that in \emph{communicating} \glspl{mdp}, the average reward $\averagereturnfunction(s)$ no longer depends on $s$.
The results of UCRL2 depend on the notion of \emph{diameter} of the \gls{mdp} which is intrinsically related to the property of being communicating. 
In fact, finite diameter and communicating \glspl{mdp} are equivalent \cite{auer2008near}. We give its definition in the following.

\begin{defi}[Diameter of an \gls{mdp}\cite{auer2008near}]
\label{def:diameter}
    Let $\markovdecisionprocess$ be an \gls{mdp} and let $\policy$ be a stationary policy on $\markovdecisionprocess$. 
    Let $T(s'\vert s;\policy)$ be the random variable that measures the time it takes for the policy $\pi$ to first reach $s'$ starting from $s$. 
    Then the diameter of $\markovdecisionprocess$ is,
    \begin{align*}
        D(\markovdecisionprocess) \coloneq \max_{s\neq s'\in\statespace}\min_{\pi\in\Pi^{\text{S}}} \expectedvalue\left[T(s'\vert s;\policy)\right].
    \end{align*}
\end{defi}

Leveraging the knowledge of the diameter, \cite{auer2008near} propose an analysis of their optimism-based algorithm that shows sub-linear regret. 
